\documentclass[aspectratio=169,11pt]{beamer}

\usepackage[utf8]{inputenc}
\usepackage{amsmath,amssymb,amsfonts}
\usepackage{physics}          % for \dv, \pdv, etc.
\usepackage{mathtools}
\usepackage{bm}
\usepackage{tikz}
\usepackage{booktabs}
\usepackage{graphicx}

\usetheme{Madrid}
\usecolortheme{default}

% Cleaner math
\setlength{\abovedisplayskip}{6pt}
\setlength{\belowdisplayskip}{6pt}
\setlength{\abovedisplayshortskip}{4pt}
\setlength{\belowdisplayshortskip}{4pt}

\title{Introduction to Plasma Physics}
\subtitle{}
\author{}
\date{}

\begin{document}

%================================================
\begin{frame}
\titlepage
\end{frame}

%================================================
\begin{frame}{What is Plasma}
\begin{columns}[c] \column{0.48\textwidth} \centering \includegraphics[width=\linewidth]{ionrecomb1.png} \\[0.3cm] \column{0.48\textwidth} \centering \includegraphics[width=\linewidth]{hyionrecomb.png} \\[0.3cm] \end{columns}

\end{frame}
\begin{frame}{Nature of Plasma}
    \begin{itemize}
        \item Quasi-neutral
        \begin{itemize}
            \item $\rho_e = e_-(n_e) + e_+(n_n) = e(n_e - n_n)$
        \end{itemize}
        \item What ever charge is dominate is quickly expelled
        \item electrostatic force $\gg$ kinetic pressure force
        \item Responds collectively to magnetic field
        \item Low energy plasma is in a constant state of Ionization, excitation, recombination, and collisions
    \end{itemize}
\end{frame}
\begin{frame}{Lorentz Force}
    \[
    F = q(\vb{E} + \vb{v} \times \vb{B})
    \]

\end{frame}




\end{document}
